\section{Kết luận và hướng phát triển}

Bài báo đã trình bày một đánh giá bán thực nghiệm đối với cơ chế Hash-linked Tamper-evident Ledger nhằm kiểm chứng tính toàn vẹn dữ liệu trong hệ thống Data Pipeline. Cơ chế này ghi nhận bằng chứng integrity ở từng stage và liên kết các ledger blocks thông qua previous hash. Verification engine tính toán lại record count, schema hash, batch hash, block hash và chain link để phát hiện modification sau xử lý.

Các quan sát thực nghiệm hỗ trợ các giả thuyết chính trong phạm vi bounded scope. Tất cả tamper scenarios được định nghĩa trước đều được phát hiện, điều kiện sạch vẫn VALID trong dashboard export quan sát được, và lineage context hỗ trợ xác định affected stage. Runtime overhead tăng khi record count tăng, nhưng kết quả overhead vẫn là approximate và chịu ảnh hưởng bởi cloud compute variability.

Nghiên cứu nên được hiểu là bằng chứng cho một cơ chế tamper-evident nhẹ, không phải bằng chứng về bảo mật Blockchain đầy đủ. Nguyên mẫu không bao gồm decentralized consensus, nhiều independent nodes, smart contracts hoặc external trust anchoring. Vì vậy, security boundary của cơ chế hiện tại bị giới hạn trong controlled workspace setting. Dù vậy, kết quả cho thấy việc kết hợp batch hash, block hash, previous hash và lineage event là một hướng khả thi để tăng khả năng kiểm chứng integrity của Data Pipeline.

Đóng góp chính của nghiên cứu là chuyển bài toán integrity từ mô tả kiến trúc sang một thiết kế thực nghiệm có câu hỏi nghiên cứu, giả thuyết, biến độc lập, biến phụ thuộc và phân tích validity. Đây cũng là điểm đáp ứng yêu cầu phương pháp luận: cơ chế kỹ thuật không được trình bày như một sản phẩm hoàn chỉnh, mà như một treatment được đánh giá trong điều kiện bán thực nghiệm.

Trong tương lai, nghiên cứu cần được lặp lại trên nhiều data platforms, mở rộng tập tamper vectors thông qua adversarial testing, tăng số lần benchmark repetitions, export raw benchmark values, externalize ledger sang một trust boundary độc lập và bổ sung digital signatures nhằm tăng cường bằng chứng chống lại privileged insider modification. Ngoài ra, có thể nghiên cứu Merkle tree hoặc partition-level hash để giảm chi phí verification khi dataset tăng lên hàng trăm nghìn hoặc hàng triệu bản ghi.
